\section{Approximate permutation fairness}\label{sec:approximate}

Exact permutation fairness forces exponentially many faces on most dice.
We now ask how small the dice can be if every order is allowed a relative
error of at most $\varepsilon$.  Thus the required inequality is
\begin{equation}\label{eq:pointwise-approx}
  \left|n!P(\pi)-1\right|\leq\varepsilon
  \qquad\text{for every }\pi\in S_n.
\end{equation}
This controls even individually rare orders.  In particular, it implies
$\lVert P-U\rVert_1\leq\varepsilon$ and
$d_{\mathrm{TV}}(P,U)\leq\varepsilon/2$, where $U$ is uniform on $S_n$.

\subsection{Classical shuffle constructions}

The periodic word $(12\cdots n)^m$ realizes an inverse $m$-riffle shuffle:
each letter independently chooses a uniform level in $[m]$, and ties are
resolved in increasing letter order.  If $d=\operatorname{des}(\pi)$,
the classical $P$-partition count gives
\begin{equation}\label{eq:descent-count}
  P(\pi)=\frac1{m^n}\binom{m+n-1-d}{n},
  \qquad
  n!P(\pi)=\prod_{k=-d}^{n-1-d}\left(1+\frac{k}{m}\right)
  \quad(m\geq n).
\end{equation}
See Stanley~\cite{stanley1972} and Bayer and Diaconis~\cite{bayer-diaconis1992}.
The product decreases with $d$.  Its two extreme values show that
$m\geq n(n-1)/\varepsilon$ suffices for \eqref{eq:pointwise-approx} when
$0<\varepsilon\leq1$: the increasing-order ratio is at most
$e^{\varepsilon/2}\leq1+\varepsilon$, and the decreasing-order ratio is at
least $1-\varepsilon/2$.  Conversely, for $0<\varepsilon<1$ the decreasing
order forces $m\geq n$ and
\[
  m\geq\frac{n(n-1)}{2[-\log(1-\varepsilon)]}.
\]
Thus this periodic family requires $\Theta(n^2/\varepsilon)$ faces per
die for $0<\varepsilon\leq1/2$.

A better classical construction repeats a palindrome:
\[
  (12\cdots n\,n\cdots21)^m.
\]
Its order distribution is exactly the shelf shuffle of Diaconis, Fulman,
and Holmes~\cite{diaconis-fulman-holmes2013}, as given by their
Description~1.  Each die independently chooses one of the \(2m\)
monotone runs uniformly; reading the chosen letters groups them by run,
in increasing order within odd runs and decreasing order within even
runs.  This is precisely their packet construction, so their
Theorem~3.1 applies with the valleys of the output permutation itself.
Their Theorem~3.4 gives, for fixed $c>0$ and $m\sim cn^{3/2}$,
\[
  \max_{\pi\in S_n}|n!P(\pi)-1|
  \longrightarrow e^{1/(12c^2)}-1.
\]
Consequently, for each fixed positive error, $O(n^{3/2})$ faces per die
suffice.  Our construction improves this bound by independently relabelling
the palindrome blocks.

\subsection{Independent palindrome blocks}

For a permutation $\rho$, let $\rho^{\mathrm{rev}}$ be its reversal.
Choose $\rho_1,\ldots,\rho_m$ independently and uniformly from $S_n$, and
form the word
\begin{equation}\label{eq:random-palindrome-word}
  W=(\rho_1\rho_1^{\mathrm{rev}})
       \concat\cdots\concat(\rho_m\rho_m^{\mathrm{rev}}).
\end{equation}
Every resulting die has exactly $2m$ faces.  The randomness here chooses
the dice; after the word has been chosen, their order distribution $P_W$
is fixed.

\begin{theorem}[Approximate fairness from random palindromes]
\label{thm:approximate}
There is an absolute constant $C$ such that, for every $n\geq3$ and
$0<\varepsilon\leq1$, if
\[
  m\geq Cn^{13/10}(\log n)^{1/5}\varepsilon^{-2/5},
\]
then the random word \eqref{eq:random-palindrome-word} satisfies
\eqref{eq:pointwise-approx} simultaneously for all $\pi\in S_n$ with
probability at least $1/2$.  In particular, there exist $n$ equal-sided
dice with
\[
  O\!\left(n^{13/10}(\log n)^{1/5}\varepsilon^{-2/5}\right)
\]
faces each satisfying \eqref{eq:pointwise-approx}.
\end{theorem}

The palindrome structure makes every pair of letters exactly fair within
each block.  Thus a block contributes no error unless at least three dice
choose faces in it.  Independent relabelling centers the remaining errors;
the proof controls their joint contribution, including interactions between
several occupied blocks.

Here is the exact expansion behind that argument.  Fix a target order
$\sigma$.  Let $L=(L_1,\ldots,L_m)$ have the multinomial distribution with
$n$ trials and $m$ equally likely cells, and split $\sigma$ into consecutive
segments $A_i$ of lengths $L_i$.  For an ordered list $A$ of distinct
letters, write $q_i(A)$ for the probability that these letters appear in
the prescribed order when each chooses one of its two occurrences in
block $i$.  Put
\[
  \delta_i(A)=|A|!q_i(A)-1,
\]
with $\delta_i(\varnothing)=0$.  Summing over the chosen block indices gives
the identity
\begin{equation}\label{eq:palindrome-expansion}
  n!P_W(\sigma)=\mathbb E_L\prod_{i=1}^m(1+\delta_i(A_i)).
\end{equation}
For $|A|\leq2$, $\delta_i(A)=0$ identically.  For every fixed $A$,
uniform relabelling gives $\mathbb E_{\rho_i}\delta_i(A)=0$.

The following estimate controls the whole expansion.  Its proof, given in
Appendix~\ref{app:approximate-proof}, combines hypercontractivity with a
bound on multinomial collision probabilities.  The $L^p$ norm is over
the independent random permutations defining $W$.

\begin{lemma}[Uniform moment estimate]\label{lem:palindrome-moment}
Let $n\geq3$, $2\leq p\leq n^2$, and
$m\geq8n^{11/10}p^{1/5}$.  For every $\sigma\in S_n$,
\[
  \bigl\|n!P_W(\sigma)-1\bigr\|_p
  \leq32\,\frac{n^{11/4}\sqrt p}{m^{5/2}}.
\]
\end{lemma}

\begin{proof}[Proof of Theorem~\ref{thm:approximate}]
Set
\[
  p=\max\left\{2,\frac{\log(2n!)}{\log3}\right\}.
\]
For $n\geq3$, we have $p\leq n^2$ and $p=O(n\log n)$.
Choosing the absolute constant $C$ sufficiently large ensures
$m\geq8n^{11/10}p^{1/5}\varepsilon^{-2/5}$.
Lemma~\ref{lem:palindrome-moment} then gives, for every $\sigma$,
\[
  \bigl\|n!P_W(\sigma)-1\bigr\|_p
  \leq\frac{32}{8^{5/2}}\varepsilon<\frac\varepsilon3.
\]
Markov's inequality and a union bound imply
\[
  \Pr\!\left(\exists\sigma\in S_n:
      |n!P_W(\sigma)-1|>\varepsilon\right)
  \leq n!3^{-p}\leq\frac12.
\]
Thus a single sampled collection satisfies all the required inequalities
with the stated probability.
\end{proof}

\subsection{A lower bound for total size}

For fixed error, quadratic total length is necessary even if only the
first-place marginals are required to be approximately fair.  Suppose
each die is first with probability at most $(1+\varepsilon)/n$.
Order the letters as $i_1,\ldots,i_n$ according to the positions
$p_1<\cdots<p_n$ of their first occurrences.  Let $E_k$ be the event
that every earlier letter $i_j$, $j<k$, chooses an occurrence after
$p_k$.  If $E_k$ fails, the first selected occurrence belongs to one
of those $k-1$ letters.  Consequently,
\[
  \Pr(E_k)\geq1-\frac{(k-1)(1+\varepsilon)}n.
\]
The choice of a face of $i_k$ is independent of $E_k$.  If $i_k$
chooses its first occurrence and $E_k$ holds, then $i_k$ is first.
Writing $m_i$ for the face counts gives
\begin{equation}\label{eq:approx-first-lower}
  m_{i_k}\geq
  \max\left\{0,\frac{n}{1+\varepsilon}-(k-1)\right\}.
\end{equation}
Summing over $k$ yields total length $\Omega(n^2)$ for fixed
$\varepsilon$.  Condition~\eqref{eq:pointwise-approx} implies the marginal
hypothesis by summing over the orders in which a specified die is first.
Theorem~\ref{thm:approximate} gives total length
$O(n^{23/10}(\log n)^{1/5})$ for fixed positive error, leaving a factor
$O(n^{3/10}(\log n)^{1/5})$ between these upper and lower bounds.
